\subsection{Contador Síncrono MOD 10}

\paragraph{Incremento}

A operação de incremento realiza a transição:
\[
D \mapsto (D + 1) \mod 10
\]


A Tabela abaixo descreve o comportamento completo da função de incremento:

\begin{center}
\begin{tabular}{c|cccc}
$D$ & $O_2$ & $O_1$ & $O_0$ & $Carry$\\
\hline
$000$ & $0$ & $0$ & $1$ & $0$\\
$001$ & $0$ & $1$ & $0$ & $0$\\
$010$ & $0$ & $1$ & $1$ & $0$\\
$011$ & $1$ & $0$ & $0$ & $0$\\
$100$ & $1$ & $0$ & $1$ & $0$\\
$101$ & $0$ & $0$ & $0$ & $1$\\
$110$ & $0$ & $0$ & $0$ & $0$\\
$111$ & $0$ & $0$ & $0$ & $0$\\
\end{tabular}
\end{center}

A partir da tabela verdade, e utilizando mapas de Karnaugh, obtêm-se as seguintes expressões booleanas:

\begin{align*}
O_2 &= (\lnot D_2 \land D_1 \land D_0)\ \lor\ (D_2 \land \lnot D_1 \land \lnot D_0) \\
O_1 &= (\lnot D_2 \land \lnot D_1 \land D_0)\ \lor\ (\lnot D_2 \land D_1 \land \lnot D_0) \\
O_0 &= (\lnot D_2 \land \lnot D_0)\ \lor\ (\lnot D_1 \land \lnot D_0) \\
Carry &= D_2 \land \lnot D_1 \land D_0
\end{align*}

O sinal $Carry$ indica a ocorrência de overflow, isto é, a transição de $101$ para $000$.

As Figuras~\ref{fig:incmod10.png} e~\ref{fig:incmod10-red.png} apresentam, respectivamente, a implementação lógica do circuito e sua realização em redstone.

\imgbig{incmod10.png}{Circuito lógico para incremento em módulo 10.}
\imgbig{incmod10-red.png}{Implementação em redstone do circuito de incremento em módulo 10.}

\paragraph{Decremento}

A operação de decremento realiza a transição:
\[
D \mapsto (D - 1) \mod 10
\]

A tabela correspondente é:

\begin{center}
\begin{tabular}{c|cccc}
$D$ & $O_2$ & $O_1$ & $O_0$ & $Carry$\\
\hline
$000$ & $1$ & $0$ & $1$ & $1$\\
$001$ & $0$ & $0$ & $0$ & $0$\\
$010$ & $0$ & $0$ & $1$ & $0$\\
$011$ & $0$ & $1$ & $0$ & $0$\\
$100$ & $0$ & $1$ & $1$ & $0$\\
$101$ & $1$ & $0$ & $0$ & $0$\\
$110$ & $0$ & $0$ & $1$ & $0$\\
$111$ & $1$ & $1$ & $0$ & $0$\\
\end{tabular}
\end{center}

As expressões booleanas minimizadas são:

\begin{align*}
O_2 &= (\lnot D_2 \land \lnot D_1 \land \lnot D_0)\ \lor\ (D_2 \land D_0) \\
O_1 &= (D_1 \land D_0)\ \lor\ (D_2 \land \lnot D_1 \land \lnot D_0) \\
O_0 &= \lnot D_0 \\
Carry &= \lnot D_2 \land \lnot D_1 \land \lnot D_0
\end{align*}

Neste caso, o sinal $Carry$ indica underflow, ou seja, a transição de $000$ para $101$.

As Figuras~\ref{fig:decmod10.png} e~\ref{fig:decmod10-red.png} mostram a implementação lógica e sua versão em redstone.

\imgbig{decmod10.png}{Circuito lógico para decremento em módulo 10.}
\imgbig{decmod10-red.png}{Implementação em redstone do circuito de decremento em módulo 10.}

\paragraph{Considerações de Projeto}

Embora funcional, o contador apresentado não é autocorretivo, pois estados inválidos não são redirecionados explicitamente para o ciclo válido. Em aplicações físicas sujeitas a ruído ou falhas transitórias, isso pode levar a comportamentos inesperados.

Uma alternativa mais robusta consiste em tratar os estados inválidos como condições de correção, forçando a transição para um estado válido (tipicamente $000$), ou utilizando tais estados como \textit{don't care} durante a minimização lógica para otimizar o circuito.
